
% TARGET cdp:st:A.1.1
\begin{proposition}[The conormal filtration and its exact section criterion]\label{cdp:st:A.1.1}
\StatementRecord{cdp:st:A.1.1}{proved}{}{the hypotheses, notation and definitions stated here; the explicit constructions and calculations in the proof}
\label{wbII:prop:criterion}\label{wbII:prop:filtration}\label{wbII:lem:conormal}\label{wbII:lem:conormal-injective}\label{wbII:prop:missing-global-section-criterion}\label{wbII:prop:missing-cotangent-filtration}
The conormal map $d:N^*\to V$, defined by $[F]\mapsto dF|_S$ for a local defining equation $F$, is injective. With $K=\ker(V\to Q)$, there are canonical exact sequences
\begin{align}
0&\longrightarrow N^*\xrightarrow{d}K\longrightarrow T\longrightarrow0,
\label{cdp:A:filtration1}\\
0&\longrightarrow K\longrightarrow V\longrightarrow Q\longrightarrow0.
\label{cdp:A:filtration2}
\end{align}
Let $\partial_A:H^0(S,Q\otimes A)\to H^1(S,K\otimes A)$ and $\delta_A:H^0(S,T\otimes A)\to H^1(S,N^*\otimes A)$ be their connecting maps. Then
\[
 H^0(S,V\otimes A)=0
 \quad\Longleftrightarrow\quad
 H^0(S,K\otimes A)=0\ \hbox{and}\ \partial_A\hbox{ is injective}.
\tag{A.1}\label{cdp:eq:A.1}
\]
Under the explicit hypotheses $H^0(S,N^*\otimes A)=0$ and $H^0(S,Q\otimes A)=0$, $H^0(S,V\otimes A)$ is canonically $\ker\delta_A$.
\end{proposition}

% TARGET cdp:st:A.4.7
\begin{proposition}[Unitary-flat differentials on every compact reduced complex surface]\label{cdp:st:A.4.7}
\StatementRecord{cdp:st:A.4.7}{proved}{}{the hypotheses, notation and definitions stated here; the explicit constructions and calculations in the proof}

Let $S$ be a compact reduced complex surface, let $\mathbb A$ be a rank-one unitary complex local system on its underlying topological space, and let $A=\mathbb A\otimes_{\C}\OO_S$ be its associated holomorphic line bundle. There is a canonical complex-linear injection
\[
 H^0(S,\widetilde\Omega_S^1\otimes A)
       \longrightarrow H^1(S,\mathbb A).
 \tag{A.20}\label{cdp:eq:A.20}
\]
The construction uses continuous local primitives that are holomorphic on the normalization; it does not assert that those primitives are holomorphic on a non-seminormal surface.
\end{proposition}

% TARGET cdp:st:A.4.9
\begin{proposition}[CDP20 Proposition 7.3 with a complete replacement proof]\label{cdp:st:A.4.9}
\StatementRecord{cdp:st:A.4.9}{proved}{}{the hypotheses, notation and definitions stated here; \cref{cdp:st:D.1.9}; \cref{cdp:st:C.3.1}; \cref{cdp:st:A.4.8}; \cite[\S1 and the central-extension diagram in the proof of Lemma~1.1]{SS2002}; \cite[\S2.1.1 and Theorem~2.2]{Albanese2020}; \cite[\S2.1.2 and Proposition~2.3]{Albanese2020}; \cref{cdp:st:D.1.8}; \cref{cdp:st:A.4.7}; \cref{cdp:st:A.4.2}; \cref{cdp:eq:A.24}}

Let $Y$ be a connected smooth compact complex threefold with $H^1(Y,\Z)=H^2(Y,\Z)=0$, algebraic dimension one and holomorphic algebraic reduction $g:Y\to C$. There is a dense open subset of $\Pic(Y)$, whose complement is a proper complex analytic subset, such that
\[
 H^0\bigl((Y_c)_{\rm red},
       \widetilde\Omega^1_{(Y_c)_{\rm red}}\otimes L|_{(Y_c)_{\rm red}}\bigr)=0
 \tag{A.23}\label{cdp:eq:A.23}
\]
for every singular reduced fibre support and every $L$ in that subset. In fact the transgression twists $L_r$ of \ref{cdp:st:D.1.9} satisfy \eqref{cdp:eq:A.23} on every reduced fibre support, including smooth and reducible supports, for every real irrational $r$.
\end{proposition}

% TARGET cdp:st:B1
\begin{proposition}[A complete nodal-family computation]\label{cdp:st:B1}
\StatementRecord{cdp:st:B1}{proved}{}{the hypotheses, notation and definitions stated here; the explicit constructions and calculations in the proof}
\label{wbII:thm:nodal-test}\label{wbII:prop:family-geometry}\label{wbII:prop:nodal-counterexample}
Let $0<\varepsilon<4/27$, let $\Delta=\{|t|<\varepsilon\}$, and define
\[
 \mathcal C=\{y^2z=x^2(x+z)+tz^3\}\subset\PP^2\times\Delta,
 \qquad X=\mathcal C\times\PP^1.
\]
The projection $X\to\Delta$ is proper, its total space is smooth, and its central fibre is $S=C_0\times\PP^1$, where $C_0$ is an irreducible nodal cubic. For each $\lambda\in\C^*$ let $A_\lambda$ be the line bundle whose pullback to the normalization $\nu:\PP^1_u\to C_0$ is trivial and whose sections in this trivialization satisfy $h(1)=\lambda h(-1)$. Put $A=\operatorname{pr}_1^*A_\lambda$ and $d_\lambda=1$ for $\lambda=1$, and $d_\lambda=0$ otherwise. In the order $q=0,1,2$, the coherent cohomology dimensions are
\[
\begin{array}{c|c}
\text{sheaf on }S & (h^0,h^1,h^2)\\\hline
 N^*_{S/X}\otimes A &(d_\lambda,d_\lambda,0)\\
 K_S\otimes A &(1,0,0)\\
 T_S\otimes A &(1,0,0)\\
 \widetilde\Omega^1_S\otimes A &(0,1+d_\lambda,d_\lambda)\\
 \Omega^1_S\otimes A &(1,1+d_\lambda,d_\lambda)\\
 \Omega^1_X|_S\otimes A &(1,1+d_\lambda,d_\lambda).
\end{array}                                                    \tag{B.1}\label{cdp:eq:B.1}
\]
Here $T_S$ denotes the torsion of the cotangent sheaf, and $K_S$ its inverse image in the ambient cotangent sheaf. In particular the two proposed vanishings hold exactly for $\lambda\ne1$, and the ambient space of sections has dimension one for every $\lambda\in\C^*$. After shrinking $\Delta$, every $A_\lambda$ extends to a line bundle on $\mathcal C$ of the form $\OO_{\mathcal C}(P-Q)$ for holomorphic sections $P,Q$.
\end{proposition}

% TARGET cdp:st:C.2.G.1
\begin{theorem}[The Euler sign for every degeneration and its equality criterion]
\label{cdp:st:C.2.G.1}
\StatementRecord{cdp:st:C.2.G.1}{proved}{}{\ExternalPaperRef{}{3.6}{II}; \ExternalPaperRef{}{5.14}{II};
\ExternalPaperRef{}{6.2}{II}; \ExternalPaperRef{}{6.3}{II};
\ref{cdp:st:C.2.MH.1}; \ref{cdp:st:C.2.MH.2};
\ExternalPaperRef{}{8.26}{II}; \ExternalPaperRef{}{8.32}{II};
\ExternalPaperRef{}{8.36}{II}; \ExternalPaperRef{}{8.37}{II};
\ref{cdp:st:C.2.INT.1}; \ref{cdp:st:C.2.INT.2};
\ref{cdp:st:C.2.INT.3}; \ref{cdp:st:C2a};
\cite[Proposition~2.1 and Theorem~2.2, pp.~842--843]{Peters2021}}

Let $Y$ be a connected smooth compact complex threefold with
$b_1(Y)=b_2(Y)=0$, and let $f:Y\to B$ be a proper surjective
holomorphic map with connected fibres to a smooth compact curve.
Assume that the smooth general fibre is a complex two-torus.
The full fibres may have arbitrary positive component multiplicities
and their reduced supports may have arbitrary singularities. Then
\[
 B\simeq\PP^1,\qquad
 c_3(Y)=e(Y)\in\{0,2\},\qquad b_3(Y)=2-c_3(Y).
 \tag{C.G.1}\label{cdp:eq:C.G.1}
\]
Every reduced fibre support has nonnegative Euler number.
An irreducible singular support with rational normalization
resolution has Euler number two; every other irreducible support
has Euler number zero.

The equality $c_3(Y)=0$ holds exactly when every support has one
of the following descriptions. A smooth irreducible support has
a finite unramified torus cover. A singular irreducible support
has smooth minimal elliptic ruled or smooth minimal Hopf
normalization. A reducible support has normalization a disjoint
union of smooth minimal Hopf or non-K\"ahler elliptic surfaces;
its reduced conductor image and inverse image are disjoint
unions of smooth elliptic curves, and the conductor branch maps
are finite unramified covers. Each description concerns the
normalization and resolution maps of the actual fibre. In
particular the full-fibre and conductor multiplicities are
retained.
\end{theorem}

% TARGET cdp:st:C.2.INT.1
\begin{proposition}[Proved elliptic-image restriction]\label{cdp:st:C.2.INT.1}
\StatementRecord{cdp:st:C.2.INT.1}{proved}{}{the hypotheses, notation and definitions stated here; \cref{cdp:eq:C.INT.3}; \ExternalPaperRef{proposition~}{3.4}{II}; \cref{cdp:eq:C.INT.1}; \cref{cdp:eq:C.INT.2}; \cite[Lemma~2.2 and \S2]{BrinzanescuMoraru2005}; \cite[Remark~2.7]{DlousskyTeleman2012}}


Let $Y$ be a smooth compact complex threefold with
$H^2(Y,\Q)=0$, let $S_i$ be an irreducible component of
a connected reducible divisor $S$, and let
$h:Z\to S_i\hookrightarrow Y$ be the minimal resolution
of its normalization. Suppose that the minimal model
$\sigma:Z\to Z_0$ is a Hopf surface or a non-K\"ahler
elliptic surface. Every strict transform $E\subset Z$ of
an elliptic curve on $Z_0$ has a one-dimensional image
under $h$, and the normalization of that image is elliptic.
The actual restriction homomorphism
$\Pic^0(Y)\to\Pic^0(E)$ is nonconstant for each such $E$.
In particular such a component cannot meet a component
whose minimal model is an Enoki surface.
\end{proposition}

% TARGET cdp:st:C.2.INT.2
\begin{proposition}[Proved exclusion of Enoki components]\label{cdp:st:C.2.INT.2}
\StatementRecord{cdp:st:C.2.INT.2}{proved}{}{the hypotheses, notation and definitions stated here; \ExternalPaperRef{proposition~}{4.3}{II}; \cref{cdp:st:C.2.INT.1}; \ExternalPaperRef{proposition~}{6.3}{II}; \cite[Lemma~2.22 and the following paragraph]{Dloussky2006}; \ExternalPaperRef{proposition~}{5.19}{II}; \cref{cdp:eq:C.INT.5}; \cref{cdp:eq:C.INT.7}; \cref{cdp:eq:C.INT.6}}


Let $f:Y\to B$ be a proper connected-fibre map from a
smooth compact complex threefold with $H^2(Y,\Q)=0$
and general fibre a complex two-torus. No component
of a reducible fibre has normalization resolution
with Enoki minimal model.
\end{proposition}

% TARGET cdp:st:C.2.INT.3
\begin{proposition}[Proved sign and equality for every reducible fibre]\label{cdp:st:C.2.INT.3}
\StatementRecord{cdp:st:C.2.INT.3}{proved}{}{the hypotheses, notation and definitions stated here; \cref{cdp:st:C.2.INT.2}; \cref{cdp:st:C.2.INT.1}; \ExternalPaperRef{proposition~}{6.3}{II}; \cref{cdp:eq:C.INT.18}; \cref{cdp:eq:C.INT.21}; \cref{cdp:eq:C.INT.20}; \cref{cdp:eq:C.INT.22}; \cref{cdp:eq:C.INT.26}; \cref{cdp:eq:C.INT.27}}


Let $f:Y\to B$ satisfy the hypotheses of \ref{cdp:st:C.2.INT.2},
and let $S$ be a connected reducible fibre support.
Then
\[
 e(S)\ge0.
 \tag{C.INT.15}\label{cdp:eq:C.INT.15}
\]
Equality holds precisely when the normalization of $S$
is a disjoint union of smooth minimal Hopf or
non-K\"ahler elliptic surfaces, and the reduced
conductor image and its reduced inverse image are
disjoint unions of smooth elliptic curves. The maps
between these elliptic curves are finite unramified
covers. This criterion allows arbitrary conductor
scheme structure and arbitrary unequal multiplicities
of the full fibre.
\end{proposition}

% TARGET cdp:st:C.2.MH.2
\begin{proposition}[Euler number from a ruled normalization]\label{cdp:st:C.2.MH.2}
\StatementRecord{cdp:st:C.2.MH.2}{proved}{}{the hypotheses, notation and definitions stated here; \ExternalPaperRef{proposition~}{3.6}{II}; \ExternalPaperRef{proposition~}{3.5}{II}; \cref{cdp:st:C.2.MH.1}; \cref{cdp:eq:C.MH.9}; \cref{cdp:eq:C.MH.8}; \cref{cdp:eq:C.MH.13}; \cref{cdp:eq:C.MH.11}; \cref{cdp:eq:C.MH.14}}

Let $f:Y\to\PP^1$ satisfy the full CDP hypotheses
$b_1(Y)=b_2(Y)=0$ with smooth torus general fibre.
Let $mS$ be a full fibre with irreducible reduced support,
and let $Z\to T\to S$ be a resolution of its
normalization. If $Z$ is rational, then
\[
                              e(S)=2.
 \tag{C.MH.6}\label{cdp:eq:C.MH.6}
\]
If $Z$ is ruled over a curve of genus $g\ge1$, then
\[
                              e(S)=2-2g.
 \tag{C.MH.7}\label{cdp:eq:C.MH.7}
\]
The calculation retains all multiplicities in the full
fibre and places no normal crossings assumption on $S$.
\end{proposition}

% TARGET cdp:st:C.2.UN.1
\begin{proposition}[Every torus-degeneration fibre has even Euler number]\label{cdp:st:C.2.UN.1}
\StatementRecord{cdp:st:C.2.UN.1}{proved}{}{the hypotheses, notation and definitions stated here; \cref{cdp:eq:C.UN.3}; \cref{cdp:eq:C.UN.4}}

Let $f:Y\to B$ be a proper holomorphic map from a smooth complex
threefold to a smooth complex curve, with connected fibres and smooth
general fibre a complex two-torus. For every $b\in B$,
\[
                    e((Y_b)_{\rm red})\in2\Z.       \tag{C.UN.1}\label{cdp:eq:C.UN.1}
\]
This assertion allows arbitrary multiplicities and singularities and
imposes no K\"ahler or Betti-number hypothesis on $Y$.
\end{proposition}

% TARGET cdp:st:C.2.UN.14
\begin{proposition}[The same non-uniruled components occur with every even Euler number]\label{cdp:st:C.2.UN.14}
\StatementRecord{cdp:st:C.2.UN.14}{proved}{}{the hypotheses, notation and definitions stated here; \cref{cdp:st:C.2.UN.4}; \cref{cdp:eq:C.EC.2}}

Fix a proper holomorphic family $p^\times:\mathcal T^\times\to\Delta^\times$
of connected complex two-tori. A marked smooth filling is a proper
holomorphic map $f:Y\to\Delta$ from a smooth complex threefold with
connected fibres, together with a biholomorphism
$f^{-1}(\Delta^\times)\simeq\mathcal T^\times$ over $\Delta^\times$.
For each irreducible component $S_i$ of $(Y_0)_{\rm red}$, let
$Z_i\to S_i^\nu\to S_i$ be a smooth compact resolution of its
normalization. Write $\mathcal E(f)$ for the finite multiset of
bimeromorphism classes $[Z_i]$ for which $Z_i$ is not uniruled.
Here a smooth compact surface is uniruled when it admits a dominant
meromorphic map $C\times\PP^1\dashrightarrow Z$ for a smooth compact
curve $C$. Repetitions in the multiset record distinct components.

For every integer $k$, the following equality holds between sets of
finite multisets:
\[
 \begin{split}
 \{\mathcal E(f): f\text{ is a marked smooth filling of }p^\times\}
 ={}&\{\mathcal E(f): f\text{ is such a filling},\\[-2pt]
 &\hspace{48mm}e((Y_0)_{\rm red})=2k\}.
 \end{split}                                                     \tag{C.EC.1}\label{cdp:eq:C.EC.1}
\]
For each filling in the left-hand set, the representative on the
right is obtained by the explicit proper holomorphic vertical
blowup map of \ref{cdp:st:C.2.UN.4}. That map induces a bijection between the
original non-uniruled components and the resulting non-uniruled
components, and a bimeromorphism of their normalization resolutions.
Every original full-fibre multiplicity is preserved by this bijection.

The same assertion holds for compact smooth total spaces over a
fixed compact base curve, with the marked family fixed away from a
finite set $Q$. One can prescribe an arbitrary even Euler number at
each point of $Q$ simultaneously, preserving the multiset and
multiplicity at each point. In particular, the collection of
non-uniruled surface classes which occur in arbitrary compact
torus-fibred fillings is exactly the collection which occur in
compact fillings all of whose reduced fibres have Euler number zero.
\end{proposition}

% TARGET cdp:st:C.2.UN.2
\begin{proposition}[All even Euler numbers and an unbounded zero-Euler family]\label{cdp:st:C.2.UN.2}
\StatementRecord{cdp:st:C.2.UN.2}{proved}{}{the hypotheses, notation and definitions stated here; \cref{cdp:st:C2b}; \cref{cdp:st:C.2.UN.1}}

The set of Euler numbers of reduced fibres among compact smooth
threefolds fibred in two-tori is exactly $2\Z$. For each $k\ge1$
there is a smooth projective such threefold with one singular reduced
normal-crossing fibre of Euler number zero consisting of a blown-up
abelian surface, $k$ ruled surfaces over a genus-two curve, and $k$
projective planes. Their incidence graph is a chain of the ruled
components attached to the abelian component, together with $k$
projective-plane leaves attached to that component.
\end{proposition}

% TARGET cdp:st:C.2.UN.4
\begin{proposition}[Every even Euler number occurs above each fixed torus filling]\label{cdp:st:C.2.UN.4}
\StatementRecord{cdp:st:C.2.UN.4}{proved}{}{the hypotheses, notation and definitions stated here; \cref{cdp:eq:C.VR.4}; \cref{cdp:eq:C.VR.6}; \cref{cdp:eq:C.VR.8}; \cref{cdp:eq:C.VR.9}; \cref{cdp:st:C.2.UN.1}}

Let $f:Y\to B$ be a proper holomorphic map from a smooth complex
threefold to a smooth complex curve, with connected fibres and smooth
general fibre a complex two-torus. Fix $b\in B$ and put
$S=(Y_b)_{\rm red}$. For every integer $k$ there is a finite composition
of blowups with smooth compact centres in successive fibres over $b$,
\[
 \pi:\widetilde Y\longrightarrow Y,\qquad
 \widetilde f=f\pi,\qquad
 e((\widetilde Y_b)_{\rm red})=2k.                 \tag{C.VR.1}\label{cdp:eq:C.VR.1}
\]
The map $\pi$ is an isomorphism over $B\setminus\{b\}$.
In particular every such filling is dominated by a smooth filling
whose central reduced support has Euler number zero. More precisely,
every further smooth model obtained by vertical blowups admits a
zero-Euler refinement by additional vertical blowups. Thus zero-Euler
models are cofinal for this domination relation, in the explicit sense
that every object is dominated by one of them.
Above every given filling there are such zero-Euler refinements with
arbitrarily many irreducible components of their central support.

An elementary operation lowering the Euler number by two consists of
one point blowup and one smooth genus-three curve blowup. On integral
cohomology its additional summands are $\Z^2$ in degrees two and four
and $\Z^6$ in degree three. In particular, whenever the Betti numbers
are finite,
\[
 (\Delta b_2,\Delta b_3,\Delta b_4)=(2,6,2),
 \qquad \Delta b_j=0\quad(j\notin\{2,3,4\}).       \tag{C.VR.2}\label{cdp:eq:C.VR.2}
\]
All the original positive full-fibre multiplicities are retained on
their strict transforms. The construction requires no normal crossings
condition on the original support.
\end{proposition}

% TARGET cdp:st:C.3.1
\begin{proposition}\label{cdp:st:C.3.1}
\StatementRecord{cdp:st:C.3.1}{proved}{}{the hypotheses, notation and definitions stated here; \cref{cdp:st:C.3.7}; \cref{cdp:st:C.3.8}; \cref{cdp:st:C.3.9}; \cref{cdp:st:C.3.10}; \cref{cdp:st:C.3.6}; \cref{cdp:st:D.2.7}; \cref{cdp:st:C.3.12}; \cite[Theorem 3.1]{Peters1990}; \cite[p.~2]{Teleman2008}; \cite[Theorem 1]{BBK2019}}

Let $f:Y\to C$ be a holomorphic algebraic reduction of a connected
compact complex threefold with $b_1(Y)=b_2(Y)=0$ and $a(Y)=1$.
The base is $C\simeq\PP^1$. Every smooth fibre $F$ has
\[
 c_1(K_F)=0\text{ in }H^2(F;\Q),\qquad
 c_2(F)=0\text{ in }H^4(F;\Z),\qquad
 \chi(F,\OO_F)=0,
\]
and contains no exceptional curve of the first kind. A very general
smooth fibre has $\kappa(F)\leq0$. The very general fibre's type belongs to the following
list: a complex two-torus, a bielliptic surface, a primary or secondary
Kodaira surface, a primary or secondary Hopf surface, or an Inoue
surface. Every listed type has Euler characteristic zero. This is a
preliminary necessary list. Statement \ref{cdp:st:C.3.7} below excludes the
bielliptic entry, and \ref{cdp:st:C.3.8} makes the very general Hopf fibre have
algebraic dimension zero. \ref{cdp:st:C.3.9} and \ref{cdp:st:C.3.10} exclude both Kodaira
entries. Statement \ref{cdp:st:C.3.6} realizes both primary and secondary Hopf
entries, and \ref{cdp:st:D.2.7} realizes a fixed abelian surface as general
fibre. Statement \ref{cdp:st:C.3.12} excludes the Inoue entry. Thus the
final very general list consists exactly of two-tori and
primary or secondary Hopf surfaces of algebraic dimension zero.
\end{proposition}

% TARGET cdp:st:C.3.10
\begin{proposition}[Proved exclusion of primary Kodaira general fibres]\label{cdp:st:C.3.10}
\StatementRecord{cdp:st:C.3.10}{proved}{}{the hypotheses, notation and definitions stated here; \cref{cdp:st:C.3.9}; \cite[pp.~373--374]{Borcea1984}; \cref{cdp:st:D.1.12}; \cref{cdp:st:D.1.9}; \cref{cdp:st:D.1.7}; \cite[\S1]{SS2002}; \cite[Remark~1]{Borcea1984}; \cref{cdp:eq:C.PK.4}; \cite[\S3]{Borcea1984}; \cref{cdp:st:D.2.3}; \cite[Theorem~2.18 and \S\S2.19--2.21]{Steenbrink1976}; \cref{cdp:eq:C.PK.5}; \cref{cdp:eq:C.PK.1}; \cref{cdp:eq:C.PK.2}; \cref{cdp:eq:C.PK.7}}


A compact complex threefold with $b_1(Y)=b_2(Y)=0$ and
$a(Y)=1$ cannot have a holomorphic algebraic reduction
$f:Y\to\PP^1$ whose general fibre is a primary Kodaira
surface. Together with \ref{cdp:st:C.3.9} this excludes both Kodaira
types using only rational Betti-number vanishing.
\end{proposition}

% TARGET cdp:st:C.3.12
\begin{proposition}[Proved exclusion of Inoue general fibres]\label{cdp:st:C.3.12}
\StatementRecord{cdp:st:C.3.12}{proved}{}{the hypotheses, notation and definitions stated here; \cref{cdp:st:C.3.1}; \cite[pp.~274, 278, 280]{Inoue1974}; \cite[Sections 2.1--2.3]{KhaledTeleman2026}; \cref{cdp:eq:C.IN.2}; \cref{cdp:eq:C.IN.3}; \cite[Theorem 3.2]{BMS2012}; \cite[Section 3.3]{BMS2012}; \cref{cdp:eq:C.IN.6}}


Under the hypotheses of \ref{cdp:st:C.3.1}, the very general smooth fibre is
not an Inoue surface. More precisely, a proper holomorphic map from
a compact connected complex manifold to a compact Riemann surface,
whose fibres over the complement of a finite set are Inoue surfaces,
has a class in $H^1(Y;\mathbb R)$ restricting nontrivially to each
such fibre. Thus these families cannot have $b_1(Y)=0$.
This assertion permits arbitrary multiplicities and arbitrary
singularities of the omitted fibres.
\end{proposition}

% TARGET cdp:st:C.3.6
\begin{proposition}[Proved actual occurrence of both Hopf types]\label{cdp:st:C.3.6}
\StatementRecord{cdp:st:C.3.6}{proved}{}{the hypotheses, notation and definitions stated here; \cref{cdp:eq:C.H.1}}


There are compact complex threefolds with $b_1=b_2=0$ and $a=1$
whose holomorphic algebraic reduction is a smooth bundle over
$\PP^1$ with primary Hopf fibre. There are also such threefolds
with secondary Hopf fibre. In these examples every fibre has Euler
number zero and $c_3=0$.
\end{proposition}

% TARGET cdp:st:C.3.7
\begin{proposition}[Proved exclusion of bielliptic fibres]\label{cdp:st:C.3.7}
\StatementRecord{cdp:st:C.3.7}{proved}{}{the hypotheses, notation and definitions stated here; \cite[Proposition~2, pp.~605--608]{Kawai1969}; \cite[Theorem~1.1]{Peters1990}}


A compact complex threefold $Y$ with $b_1(Y)=b_2(Y)=0$ has no
proper connected-fibre holomorphic map to a compact curve whose
general fibre is bielliptic.
\end{proposition}

% TARGET cdp:st:C.3.8
\begin{proposition}[Proved exclusion of elliptic Hopf general fibres]\label{cdp:st:C.3.8}
\StatementRecord{cdp:st:C.3.8}{proved}{}{the hypotheses, notation and definitions stated here; \cref{cdp:st:C.3.1}; \cite[Proposition~3.1 and its proof]{SS2002}}


The very general Hopf fibre of the algebraic reduction in \ref{cdp:st:C.3.1}
has algebraic dimension zero. In particular neither primary nor
secondary Hopf surfaces of algebraic dimension one can be its
general fibre. This exclusion uses no restriction-surjectivity
assertion from CDP98 or CDP20.
\end{proposition}

% TARGET cdp:st:C.3.9
\begin{proposition}[Proved exclusion of secondary Kodaira general fibres]\label{cdp:st:C.3.9}
\StatementRecord{cdp:st:C.3.9}{proved}{}{the hypotheses, notation and definitions stated here; \cref{cdp:st:D.1.9}; \CompendiumRef{II}{cdp:st:E4}; \cref{cdp:st:D.1.7}; \cref{cdp:eq:C.SK.1}; \cite[Theorem~4 and Remark~1]{Borcea1984}; \cite[\S1, pp.~2--3]{SS2002}; \cref{cdp:eq:C.SK.2}; \cref{cdp:eq:C.SK.3}; \cref{cdp:eq:C.SK.4}; \cref{cdp:eq:C.SK.5}}


A compact complex threefold with $b_1(Y)=b_2(Y)=0$ and
$a(Y)=1$ cannot have a holomorphic algebraic reduction
$f:Y\to\PP^1$ whose general fibre is a secondary Kodaira
surface. Rational Betti-number vanishing suffices; no integral
torsion-freeness of the fibre is assumed.
\end{proposition}

% TARGET cdp:st:C4f
\begin{proposition}[An unrestricted geometric replacement for CDP20 Theorem 2.2]\label{cdp:st:C4f}
\StatementRecord{cdp:st:C4f}{proved}{}{the hypotheses, notation and definitions stated here; \cref{cdp:st:C4e}; \cref{cdp:st:C4c}; \cref{cdp:st:A.3.4}; \cref{cdp:st:A.3.5}; \cref{cdp:st:A.5.2}; \cref{cdp:st:A.2.2}; \cref{cdp:eq:C.VC.4}; \cref{cdp:st:A.5.13}; \cref{cdp:st:C4d}; \cref{cdp:st:D.1.11}; \cref{cdp:st:A.4.9}}

Let $f:Y\to\PP^1$ satisfy the integral CDP hypotheses.
For each critical full fibre with
irreducible nonnormal reduced support $S$, let
$\eta_S:T_S\to S$ be its normalization. Then a general
line bundle $L\in\Pic(Y)$ satisfies
\[
 H^q(Y,T_Y\otimes L)=0\quad(q\ne1)
 \tag{C.VC.1}\label{cdp:eq:C.VC.1}
\]
exactly when, for every such $S$, the holomorphic pullback
homomorphism
\[
 \Pic(Y)\longrightarrow\Pic(T_S)
 \tag{C.VC.2}\label{cdp:eq:C.VC.2}
\]
is nonzero. Equivalently, its differential
$H^1(Y,\OO_Y)\to H^1(T_S,\OO_{T_S})$ is nonzero.
In this case $\chi(Y,T_Y\otimes L)\le0$ and $c_3(Y)\le0$.
If one map \eqref{cdp:eq:C.VC.2} is zero, then
$H^2(Y,T_Y\otimes L)\ne0$ for every $L\in\Pic(Y)$.
All fibre multiplicities and all reduced-support
singularities are permitted in this criterion.
\end{proposition}

% TARGET cdp:st:C5b
\begin{proposition}[Integral twists and the normalization-map exception]\label{cdp:st:C5b}
\StatementRecord{cdp:st:C5b}{proved}{}{the hypotheses, notation and definitions stated here; \cref{cdp:eq:C.8}; \cref{cdp:st:C5a}; \cref{cdp:eq:C.9}}

For $n\in\Z$ define
\[
 r_1(n)\equiv n\pmod3,\quad 0\le r_1(n)<3,
 \qquad r_2(n)\equiv3n\pmod4,\quad0\le r_2(n)<4,
\]
\begin{gather*}
 d(n)=\frac{n-4r_1(n)-3r_2(n)}{12},\qquad
 \epsilon_n=\begin{cases}0&3\mid n,\\1&3\nmid n.\end{cases}
\\
 u(n)=\max\{d(n)+1,0\},\qquad
 v(n)=\max\{-d(n)-1,0\}.
\end{gather*}
The coarse direct images of $T_X\otimes M_n$ are
\begin{align*}
 f_*(T_X\otimes M_n)&=\OO(d(n))\oplus\OO(d(n-1)),\\
 R^1f_*(T_X\otimes M_n)&=\OO(d(n))\oplus\OO(d(n-1))\oplus\OO(d(n+2)),\\
 R^2f_*(T_X\otimes M_n)&=\OO(d(n+6))\oplus\C_{p_1}^{\epsilon_n}.\tag{C.5}\label{cdp:eq:C.5}
\end{align*}
Their global cohomology dimensions are
\begin{align*}
 h^0&=u(n)+u(n-1),\\
 h^1&=v(n)+v(n-1)+u(n)+u(n-1)+u(n+2),\\
 h^2&=v(n)+v(n-1)+v(n+2)+u(n+6)+\epsilon_n,\\
 h^3&=v(n+6).                                             \tag{C.6}\label{cdp:eq:C.6}
\end{align*}
The natural map $H^2(X,T_X\otimes M_n)\to H^2(\mathrm dP_6,K_{\mathrm dP_6})=\C$ has rank one exactly for
\[
 n\in\{-6,-3,-2\}\cup\Z_{\ge0};                         \tag{C.7}\label{cdp:eq:C.7}
\]
it is zero for all other integers. In particular $M_{-1}$ has tangent cohomology $(0,0,1,0)$ and a zero map to the normalization target.
\end{proposition}

% TARGET cdp:st:C5c
\begin{proposition}[Every nonintegral twist and its resonance lengths]\label{cdp:st:C5c}
\StatementRecord{cdp:st:C5c}{proved}{}{the hypotheses, notation and definitions stated here; \cref{cdp:eq:C.13}}

Retain the fixed marked functions $\tau,\mu,\beta_c$ and put $\rho=\tau(z_1)$. For $z\notin\Z$, let $\mathcal R_z$ be the set of regular base points at which integers $(n,p,q,r)$ satisfy
\[
 q\tau+r\mu=p,\qquad z=n-6q\mu-r\beta_c.               \tag{C.10}\label{cdp:eq:C.10}
\]
Let
\[
 \Lambda_1=\Z+(3\beta_c(z_1)+4-2\rho)\Z,\quad
 \Lambda_2=\Z+(2\beta_c(z_2)+3-3i)\Z.
\]
Set $\ell_1(z)=1+\epsilon_n$ for $z=n+k(3\beta_c(z_1)+4-2\rho)$ with $k\ne0$, and zero off $\Lambda_1$. Set $\ell_2(z)=1$ for $z=n+k(2\beta_c(z_2)+3-3i)$ with $k\ne0$ and $n\equiv1,2\pmod4$, and zero otherwise. Then $\mathcal R_z$ is finite. With $N_z=|\mathcal R_z|$ and $J_z=2N_z+\ell_1(z)+\ell_2(z)$, we have
\[
 (h^0,h^1,h^2,h^3)(X,T_X\otimes M_z)=(0,J_z,1+J_z,0),  \tag{C.11}\label{cdp:eq:C.11}
\]
and, writing $\mathfrak m_t$ for the maximal ideal of $\OO_{B,t}$,
\begin{align*}
R^1f_*(T_X\otimes M_z)&=
 \bigoplus_{t\in\mathcal R_z}\C_t^2
 \oplus\C_{p_1}^{\ell_1(z)}\oplus\C_{p_2}^{\ell_2(z)},\\
R^2f_*(T_X\otimes M_z)&=
 \C_{p_0}\oplus\bigoplus_{t\in\mathcal R_z}\OO_{B,t}/\mathfrak m_t^2
 \oplus\C_{p_1}^{\ell_1(z)}\oplus\C_{p_2}^{\ell_2(z)}.\tag{C.12}\label{cdp:eq:C.12}
\end{align*}
The natural map to $H^2(\mathrm dP_6,K_{\mathrm dP_6})$ is surjective for all these nonintegral twists.
\end{proposition}

% TARGET cdp:st:D.1.2
\begin{proposition}\label{cdp:st:D.1.2}
\StatementRecord{cdp:st:D.1.2}{proved}{}{the hypotheses, notation and definitions stated here; \CompendiumRef{II}{cdp:st:D.1.1}; \cref{cdp:eq:D1}}

Let $f:Y\to C$ be a proper holomorphic surjection from a connected
compact complex threefold to a compact Riemann surface, with
connected fibres. Let $\Delta$ be a nonempty finite set containing
the critical values, $s=|\Delta|$, $g=g(C)$,
$D=f^{-1}(\Delta)_{\rm red}$, and let $r$ be the number of
irreducible components of $D$. Write $F$ for a smooth fibre and
\[
t'=\dim_{\Q}H_1(F;\Q)_{\pi_1(C\setminus\Delta)}
   =\dim_{\Q}H^1(F;\Q)^{\pi_1(C\setminus\Delta)}.
\]
For $b_2(Y)=0$ the exact component count is
\[
\boxed{\quad r=2g+s-1+t'-b_1(Y).\quad}
\]
Under CDP20's integral assumptions
$H^1(Y;\Z)=H^2(Y;\Z)=0$, the base is $\PP^1$ and
\[
r=s-1+t',\qquad
H_1(F;\Z)_{\pi_1(C\setminus\Delta)}\text{ is free abelian}.
\]
The last assertion concerns the indicated quotient of $H_1(F;\Z)$;
it does not assert that $H_1(F;\Z)$ itself is torsion-free.
\end{proposition}

% TARGET cdp:st:D.2.3
\begin{proposition}[Proved logarithmic period obstruction]\label{cdp:st:D.2.3}
\StatementRecord{cdp:st:D.2.3}{proved}{}{the hypotheses, notation and definitions stated here; \cite[\S2.1, Theorem~2.18, \S2.19, Proposition~2.20 and
Theorem~2.21]{Steenbrink1976}; \cite[Proposition~1.13 and proof of
Proposition~2.20]{Steenbrink1976}; \cref{cdp:eq:D.18}; \cref{cdp:eq:D.19}}


Let $h:U\to\Delta$ be a proper holomorphic map from a smooth
complex threefold, with connected fibres, smooth two-torus fibres
over $\Delta^*$, and central fibre $F=\sum_i m_iD_i$ with
normal-crossings reduced support $D=\bigcup_iD_i$.
If every normalization component of $D$ has geometric genus zero,
the monodromy on $H^2(U_t;\C)$ is nontrivial. In particular the
monodromy on $H^1(U_t;\Q)$ is nontrivial. No K\"ahler assumption
on $U$ or $D$ is made, and the positive integers $m_i$ are arbitrary.
\end{proposition}

% TARGET cdp:st:D.2.5
\begin{proposition}[Proved unrestricted sign under trivial monodromy]\label{cdp:st:D.2.5}
\StatementRecord{cdp:st:D.2.5}{proved}{}{the hypotheses, notation and definitions stated here; \cref{cdp:st:C4c}; \cref{cdp:st:A.3.5}; \cref{cdp:eq:D.S.2}; \cref{cdp:st:D.1.9}; \cref{cdp:st:A.4.7}; \cref{cdp:st:A.5.2}; \cref{cdp:eq:D.S.4}; \cref{cdp:st:A.5.6}; \cref{cdp:st:A.5.3}; \cref{cdp:eq:D.S.5}; \cite[Theorem~1.6]{BierstoneMilman1997}; \cref{cdp:st:D.2.3}}


Let $f:Y\to\PP^1$ be a proper surjective holomorphic map with
connected fibres from a compact connected complex threefold with
$b_1(Y)=b_2(Y)=0$. Suppose its smooth general fibre is a complex
two-torus and its monodromy representation on $H_1(F,\Q)$ is
trivial. Then
\[
 c_3(Y)\le0,\qquad b_3(Y)\ge2.
 \tag{D.S.1}\label{cdp:eq:D.S.1}
\]
The fibres may have arbitrary singularities and arbitrary
multiplicities. In particular this proves the sign assertion of
D.2 for a holomorphic algebraic reduction; no existence assertion
is required for the conclusion.
\end{proposition}

% TARGET cdp:st:D.2.7
\begin{proposition}[Proved existence]\label{cdp:st:D.2.7}
\StatementRecord{cdp:st:D.2.7}{proved}{}{the hypotheses, notation and definitions stated here; \cref{cdp:st:D.2.5}; \cref{cdp:eq:D.24}; \cref{cdp:eq:D.25}; \cref{cdp:eq:D.26}; \cref{cdp:eq:D.23}; \cref{cdp:eq:D.28}; \cref{cdp:eq:D.30}; \cref{cdp:eq:D.32}; \cref{cdp:eq:D.31}}


There exists a smooth compact connected complex threefold $Y$
with a holomorphic map $f:Y\to\PP^1$ with connected fibres such
that
\[
 H_1(Y;\Z)=H_2(Y;\Z)=0,\qquad
 a(Y)=1,\qquad c_3(Y)=0,\qquad b_3(Y)=2.
\tag{D.21}\label{cdp:eq:D.21}
\]
Every smooth fibre is the fixed abelian surface
$A=E\times E$, where $E=\C/(\Z+\mathrm i\Z)$, and the marked
integral monodromy on $H_1(A;\Z)$ is the identity. There are
exactly three singular fibres. Each is reduced with simple
normal crossings, and consists of a primary Kodaira surface
and a primary Hopf surface meeting along one smooth elliptic
curve. Consequently the trivial-monodromy class in D.2 is
nonempty, and equality in the unrestricted sign theorem
\ref{cdp:st:D.2.5} is attained.
\end{proposition}

% TARGET cdp:st:E12
\begin{proposition}[Counterexample: The universal local-freeness assertion of CDP98, Lemma 1.5]\label{cdp:st:E12}
\StatementRecord{cdp:st:E12}{proved}{}{the hypotheses, notation and definitions stated here; the explicit constructions and calculations in the proof}

Let $A$ be an elliptic curve with origin $0$, let
$P=\mathcal O_{A\times A}(\Delta-A\times\{0\}-\{0\}\times A)$,
and let $Y=A\times A\times\mathbb P^1$, with $g:Y\to A$
projection to the second factor. For
$V=P\boxtimes\mathcal O_{\mathbb P^1}(-2)$ one has
\[
 g_*V=R^1g_*V=0,\qquad R^2g_*V=\mathbb C_0,
 \qquad h^\bullet(Y,V)=(0,0,1,0).
\]
The sheaf in degree two has length one. For $t\ne0$ the fibre
cohomology is zero in every degree; for $t=0$ it is $(0,1,1)$.
\end{proposition}

% TARGET cdp:st:E13
\begin{proposition}[The Kuranishi germ of the two-parameter family]\label{cdp:st:E13}
\StatementRecord{cdp:st:E13}{proved}{}{the hypotheses, notation and definitions stated here; \cite[Proposition labelled \texttt{family}]{Viaclovsky2026Two}; \cref{cdp:eq:E13-connection}; \cref{cdp:eq:E13-residue}; \cref{cdp:st:E11}; \cref{cdp:st:E1}; \cref{cdp:st:E6}}

For every member of the family of
\cite[Proposition labelled \texttt{family}]{Viaclovsky2026Two}, on the parameter
bidiscs with $\mu$ sufficiently close to zero and nonzero sufficiently
small linearization scale $\lambda$, one has
\[
 (h^0,h^1,h^2,h^3)(X_{\mu,\lambda},T_X)=(1,2,2,0).
\]
Its Kuranishi germ is a smooth reduced surface, and the displayed
$(\mu,\lambda)$-family is locally miniversal. In particular its
sufficiently small deformations all have algebraic dimension one.
\end{proposition}

% TARGET cdp:st:E14
\begin{proposition}[The Liu--Rollenske family and its transverse deformation]\label{cdp:st:E14}
\StatementRecord{cdp:st:E14}{proved}{}{the hypotheses, notation and definitions stated here; \cite[Proposition labelled
\texttt{prop: surface facts}]{LiuRollenske2026}; \cref{cdp:eq:E14-elliptic-family}; \cite[Theorem labelled \texttt{thm: main}]{LiuRollenske2026}; \cref{cdp:st:E13}; \CompendiumRef{II}{cdp:st:E4}; \cref{cdp:eq:E14-canonical}; \cref{cdp:eq:E14-jet-map}; \cref{cdp:eq:E14-log-coho}; \cref{cdp:st:E6}; \cref{cdp:st:E1}}

Every member $Y(u)$ of the Liu--Rollenske family, with its scale
sufficiently small and nonzero as in its construction, satisfies
\[
 h^\bullet(Y(u),T_{Y(u)})=(1,2,2,0).
\]
Its Kuranishi germ is a smooth reduced surface. One parameter is
the given scale; a second parameter replaces its cuspidal fibre by
two rank-one nodal fibres. All sufficiently small deformations
retain a torus fibration and have algebraic dimension one.
\end{proposition}

% TARGET cdp:st:E15
\begin{proposition}[Tangent cohomology and Kuranishi germs of positive combs]\label{cdp:st:E15}
\StatementRecord{cdp:st:E15}{proved}{}{the hypotheses, notation and definitions stated here; \cite{Viaclovsky2026Combs}; \cref{cdp:eq:E15-products}; \cref{cdp:eq:E15-positive}; \cref{cdp:st:E5}; \cref{cdp:st:E13}; \cref{cdp:st:E1}; \cref{cdp:st:E6}}

For each positive-comb manifold $X_m$ constructed in
\cite{Viaclovsky2026Combs}, in particular for every $m\geq12$ in
its stated construction, one has
\[
 h^\bullet(X_m,T_{X_m})=(2m+1,\,8m+1,\,6m+1,\,0).
\]
Its Kuranishi germ is smooth and reduced of dimension $8m+1$.
A locally miniversal family is supplied by all sufficiently nearby
admissible period sections on the root curve, with the affine
attachment classes fixed. Its members have algebraic dimension one
and torus fibrations with the same three critical values,
multiplicities $3,4$, and hexagonal fibre type.
\end{proposition}

% TARGET cdp:st:E17
\begin{proposition}[Counterexample: The embedded-component omission in the divisor-twisting lemma]\label{cdp:st:E17}
\StatementRecord{cdp:st:E17}{proved}{}{the hypotheses, notation and definitions stated here; \CompendiumRef{II}{cdp:st:E2}}

CDP98, Proposition~1.1, in its printed generality of arbitrary
nonreduced compact spaces, is false. The exact replacement is as
follows. For an effective Cartier divisor $D$ on a compact complex
space $Z$ and a locally free sheaf $\mathcal V$, the spaces
$H^0(Z,\mathcal V(-kD))$, viewed inside $H^0(Z,\mathcal V)$,
stabilize for $k\gg0$ at
\[
 \{s\in H^0(Z,\mathcal V):\operatorname{supp}(s)\cap|D|
       =\varnothing\}.
\]
In particular they eventually vanish when $Z$ has no embedded
associated components and $D$ meets every irreducible component.
This last case applies to every full Cartier fibre used in \CompendiumRef{II}{cdp:st:E2}.
\end{proposition}

% TARGET cdp:st:E18
\begin{proposition}[Counterexample: The square theta formula for arbitrary sections]\label{cdp:st:E18}
\StatementRecord{cdp:st:E18}{proved}{}{the hypotheses, notation and definitions stated here; \cref{cdp:eq:E-theta-general}; \cref{cdp:st:E5}}

The square-only formula in Lead~9 is false as a local period
identity for arbitrary elliptic sections and linearizations.
The exact formula for the theta expression specified there is
\eqref{cdp:eq:E-theta-general}. The assertion for the actual original
period family remains true, by its equation $\wp(\mu)=0$.
\end{proposition}

% TARGET cdp:st:E5
\begin{proposition}[The infinitesimal map and the rank-two boundary]\label{cdp:st:E5}
\StatementRecord{cdp:st:E5}{proved}{}{the hypotheses, notation and definitions stated here; \CompendiumRef{II}{cdp:st:E4}}

Let $g:Y\to C$ be a proper surjective holomorphic map with connected fibres
from a smooth connected compact complex threefold to a nonsingular compact
connected complex curve. Its smooth general fibre is a two-dimensional complex
torus. There is a canonical exact sequence of vector spaces
\begin{equation}\label{cdp:eq:E-Leray-map}
0\to H^1(C,T_C)\xrightarrow{g^*}H^1(Y,g^*T_C)
 \to H^0(C,T_C\otimes R^1g_*\mathcal O_Y)\to0.
\end{equation}
On the smooth locus $U\subset C$, the relative tangent sequence
induces a sheaf morphism
\[
 \delta:T_C\otimes R^1g_*\mathcal O_Y|_U
                  \longrightarrow R^2g_*T_{Y/C}|_U.
\]
At a smooth torus fibre $F$, write its Kodaira--Spencer tensor as
$\kappa=u_1\otimes\eta_1+u_2\otimes\eta_2$, where
$\eta_1,\eta_2$ is a basis of $H^1(F,\mathcal O_F)$ and
$u_i\in\operatorname{Lie}(F)$. Then
\[
 \delta(a\eta_1+b\eta_2)
   =(bu_1-au_2)\otimes(\eta_1\wedge\eta_2).
\]
Its rank is the rank of $\kappa$. When this rank is two at one
regular value, the image of $H^1(Y,T_Y)\to H^1(Y,g^*T_C)$ lies
in $g^*H^1(C,T_C)$. For $C=\mathbb P^1$ this image is zero.
\end{proposition}

% TARGET cdp:st:E6
\begin{proposition}[Stability of a rank-two torus fibration]\label{cdp:st:E6}
\StatementRecord{cdp:st:E6}{proved}{}{the hypotheses, notation and definitions stated here; \cref{cdp:st:E5}; \cite[\S4, equations (4.3)--(4.5), and \S6]{Horikawa1974}; \cref{cdp:eq:E-surj-embedding}; \cref{cdp:eq:E-Leray-map}; \cite[Section~9, Theorems~1--2, pp.~83--85]{Douady1966}; \cite[Section~10.1, Proposition~1, p.~86]{Douady1966}; \cite[Section~10.2, Theorem~1, p.~87]{Douady1966}; \cite[Theorem 11]{Buchdahl1999}; \cref{cdp:st:E1}}

Let $g:Y\to\mathbb P^1$ be as in \ref{cdp:st:E5} and suppose its
Kodaira--Spencer tensor has rank two at one regular value. For every
proper deformation $\mathcal Y\to(S,0)$ of $Y$ over a complex
analytic germ, including a nonreduced germ, there is, after shrinking
$S$, a holomorphic map
\[
 G:\mathcal Y\longrightarrow\mathbb P^1\times S
\]
over $S$ whose central restriction is $g$. Every $G_s$ is surjective
with connected fibres and smooth general fibre a complex two-torus.
With $b_2(Y)=0$, every sufficiently small deformation has algebraic
dimension one.
\end{proposition}

% TARGET cdp:st:E8
\begin{proposition}[The theta determinant and its exact rank locus for $X_c$]\label{cdp:st:E8}
\StatementRecord{cdp:st:E8}{proved}{}{the hypotheses, notation and definitions stated here; the explicit constructions and calculations in the proof}

Use the period convention
$\Omega=\left(\begin{smallmatrix}\tau&\mu\\\mu&\beta/6\end{smallmatrix}\right)$,
Jacobi's odd theta function with periods $1,\tau$, and
$\beta/6=(i/\pi)\log(\vartheta_1(\mu,\tau)/\eta(\tau))+\text{constant}$.
Writing $h=\partial_\mu\log\vartheta_1$, one has the identity
\begin{equation}\label{cdp:eq:E-theta-general}
 \det\frac{d\Omega}{d\tau}
 =-\left(\mu'-\frac{i}{2\pi}h\right)^2
       -\frac{\wp(\mu,\tau)}{4\pi^2}.
\end{equation}
For the actual AC period branch $\wp(\mu,\tau)=0$, this becomes
\[
 \det\frac{d\Omega}{d\tau}
       =-\frac{g_2(\tau)^2}{36\pi^2g_3(\tau)}.
\]
On the smooth root-family base, the Kodaira--Spencer tensor has rank
one at the order-three point and rank two at every other finite
point, including the order-four point. Its determinant has order
two at the order-three point. The logarithmic tensor at the cusp
has rank two.
\end{proposition}

% TARGET ii:LR:covering
\begin{theorem}[The retained covering-parameter result]\label{ii:LR:covering}
\StatementRecord{ii:LR:covering}{proved}{}{\ref{cdp:st:E1}; Paper~IV, \ExternalPaperRef{}{3.33}{IV} and \ExternalPaperRef{}{3.39}{IV}; the retained constructive maps of \ExternalPaperRef{}{8.17}{IV} identified by \ExternalPaperRef{}{6.7}{IV}; \cite[Lemma~2.7, Corollary~2.8, Proposition~2.12, Proposition~3.6 and Definition~3.7]{LiuRollenske2026}}
Retain the exact scalar $u$ in the Liu--Rollenske action $B_u=uB_0$,
the good fourth root at the $III^*$ source fibre, the specified
logarithmic vector $a=e_4/4$, and its fixed line character.
The resulting construction is a holomorphic family on the universal
logarithmic cover of a sufficiently small punctured $u$-disc.
The threefold reached after one positive $u$-turn is not
biholomorphic to the initial threefold. The actual family descends
to the degree-two parameter cover
\[
                         u=w^2.\tag{LR.1}
\]
This degree is minimal. The conclusion holds with every retained
constant root-line character.
\end{theorem}

% TARGET inv:E2.4
\begin{proposition}[Counterexample: The universal two-vector-field bound]\label{inv:E2.4}
\StatementRecord{inv:E2.4}{proved}{}{the hypotheses, notation and definitions stated here; \citet[Corollary 1.3; verified preprint p.~2]{HKP2000}; \citet[Remarks 1.2 and 9.26]{AC2026}; \cref{inv:E2.2}}

The assertion that every complex structure on $S^6$ has
$h^0(T)\le2$, printed in
\citet[Corollary 1.3; verified preprint p.~2]{HKP2000}
and repeated in \citet[Remarks 1.2 and 9.26]{AC2026}, is false
for the comb members.
\end{proposition}

% TARGET night:P3K.36
\begin{refuted}[The historical all-fan clause of P3K.36]\label{night:P3K.36}\label{stmt:P3K.36}
\StatementRecord{night:P3K.36}{false as stated}{}{Paper~IV, \ExternalPaperRef{}{6.6}{IV} and \ExternalPaperRef{}{6.7}{IV}}
\par\noindent\textbf{Refined receiving item:} \ExternalPaperRef{proposition~}{8.18}{IV}. The historical compound statement and argument remain here at their recorded scope.\par

The earlier clause assigned the value
$3(S-1)(S-3)/8\pmod {28}$ to every fan of the corrected
$u=w^2$ family, where $S=p^2-pq+q^2$ was the norm of the
physical delta cocharacter. The actual fan $T_2$ has $S=1$,
so that expression gives zero. The retained degree-four
page has $(H,\chi)=(2,1)$ and its degree-two quotient has
$(J,\widehat H,\widehat\chi)=(2,1,1/2)$.
The actual fan maps in Paper~IV, \ExternalPaperRef{}{6.6}{IV} and
\ExternalPaperRef{}{6.7}{IV}, give
\[
 (\widehat\chi-\widehat H/2)-J\,2(2-1)/2
                           =-2=26\pmod {28}.
\]
This is the stated counterexample to the all-fan clause.
The original fixed reference fan retains its separately
proved zero value; the single $u$-turn has no endpoint
identification by \ref{ii:LR:covering}.
\end{refuted}

% TARGET night:P4.58
\begin{proposition}[P4.58]\label{night:P4.58}\label{stmt:P4.58}
\StatementRecord{night:P4.58}{proved}{}{\ref{night:P4.10}; \ref{night:P4.11}; \ref{night:P4.26}; \ref{night:P4.32}; \ref{night:P4.35}; \ref{night:P4.36}; \ref{night:P4.37}; \ref{night:P4.54}; \ref{night:P4.55}; \ref{night:P4.56}; \ref{night:P4.57}; \ref{night:P4.6}; \ref{night:P4.7}; \ref{night:P4.8}; \ref{night:P4R.13}; \ref{night:P4R.18}; \ref{night:P4R.20}}

Let $f:Y\to B$ be the fibration of Part~4.2, with
$b_1(Y)=b_2(Y)=0$, connected fibres and general fibre a
complex two-torus. Suppose every reduced fibre support
is irreducible or has normal crossings. Then
\[
 B=\mathbb P^1,\qquad c_3(Y)\in\{0,2\},
                       \qquad b_3(Y)=2-c_3(Y).
\]
Every reduced support has nonnegative Euler number.
An irreducible singular support whose normalization has
a rational resolution has Euler number two. Every other
irreducible support has Euler number zero. Such a support
is either smooth with a finite unramified torus cover,
or singular with smooth minimal elliptic ruled
normalization, or singular with smooth minimal Hopf
normalization. In particular the class-VII normalization
alternative is realized only by a smooth minimal Hopf
surface.

For a normal-crossing support the Euler number equals
its number of triple points, and is even. Its zero case
is exactly P4.8: in the reducible case the normalization
is a disjoint union of smooth minimal Hopf or non-K\"ahler
elliptic surfaces, with smooth elliptic conductor curves
and finite unramified branch maps. The equality
$c_3(Y)=0$ holds exactly when every support has one of
the corresponding zero-Euler forms.
\end{proposition}

% TARGET push:src:E.2.23
\begin{theorem}[separation under arbitrary deformation chains]\label{push:src:E.2.23}
\StatementRecord{push:src:E.2.23}{proved}{}{the hypotheses, notation and definitions stated here; \cref{push:src:E.2.4}; \cref{push:src:E.2.22}; \citet[Chapter~II, Theorem~4.19, pp.~95--96]{Demailly2012E}; \cite[Chapter~II, Theorems~7.6 and~7.12]{Demailly2012E}}

In a proper smooth holomorphic family of compact complex
manifolds over a connected reduced complex analytic base,
if one fibre is a sphere member of \ExternalPaperRef{theorem~}{20.6}{I} with coprime
multiplicities $(3k,4l)$, then every fibre is a constructed
member with that same pair.

Consequently two of the constructed sphere members are
deformation equivalent, allowing chains of proper smooth
families over connected complex analytic bases, if and
only if their ordered multiplicity pairs agree.  In
particular the distinct coprime pairs give distinct
deformation classes, and \ref{push:src:E.2.4} is resolved by this theorem.
\end{theorem}

% TARGET push:src:E.2.24
\begin{corollary}[separation for continuous complex structures]\label{push:src:E.2.24}
\StatementRecord{push:src:E.2.24}{proved}{}{the hypotheses, notation and definitions stated here; \citet[Section~3.2, Theorem~3.4 and Remark~3.6,
pp.~898--899]{Meersseman2019E}; \cite[Section~2.1, pp.~883--884]{Meersseman2019E}; \cref{push:src:E.2.23}}

Let $\mathcal I(S^6)$ be the space of smooth integrable
almost-complex structures on the fixed smooth sphere,
with its $C^\infty$ topology.  For a fixed coprime pair
$(3k,4l)$, let $\mathcal I_{k,l}$ consist of the structures
whose complex manifolds are biholomorphic to constructed
members with that pair.  Then $\mathcal I_{k,l}$ is open
and closed in $\mathcal I(S^6)$.

Consequently a continuous family of such structures
over any connected topological parameter space cannot
contain constructed members with two different pairs.
The same separation holds in the quotient spaces by
$\operatorname{Diff}_0(S^6)$ or by all orientation-preserving
diffeomorphisms, with their quotient topologies.
This statement does not identify the number of marked
connected components belonging to a single pair.
\end{corollary}